

Now, we are going to study the relationship between the $\mathbb{Q}$-linear series and moment polytopes.

For $m \in \mathbb{N}(L) \setminus \{0\}$, we have the map $\div : H^0(X, m L) \setminus \{0\} \surj \ls{m L}$; by composing it with $\ls{m L} \to {1 \over m} \ls{m L} \subseteq \Qls{L}$, we get a normalized version of $\div$ as
\begin{equation}
\begin{aligned}
	\divQ : H^0(X, m L) \setminus \{0\} &\to \Qls{L} , \\
	s &\mapsto {1 \over m} \div(s) .
\end{aligned}
\end{equation}

After fixing a $G$-linearization $\alpha$, for every $m \in \mathbb{N}(L, \alpha)$, space $H^0(X, m L)$ is a non-zero $G$-module, and we can restrict $\divQ$ to
\begin{equation*}
	\divQ: H^0(X, m L)^{(B)} \setminus \{0\} \surj {1 \over m} \ls{m L}^B .
\end{equation*}


We can define the following map,
\begin{equation}
\begin{aligned}
	p_\alpha: \Qls{L}^B &\to P(X, L, \alpha), \\
	D &\mapsto \lambda / m,
\end{aligned}
\end{equation}
where $D = \divQ(s)$ for some $m \in \mathbb{N}(L, \alpha)$, $\lambda \in \ch(B)$, and $s \in H^0(X, m L)^{(B)}_\lambda$.

\begin{proposition}
	The map $p_\alpha$ is a well-defined $\mathbb{Q}$-convex-linear surjective map. If the variety $X$ is spherical, the map $p_\alpha$ is a convex-linear isomorphism.
\end{proposition}
\begin{proof}
	First, to show it is well-defined (then surjectivity is from the definition of moment polytopes), for every two sections
	\begin{equation*}
		s \in H^0(X, m L)^{(B)}_\lambda, \text{ and } s' \in H^0(X, m' L)^{(B)}_{\lambda'},
	\end{equation*}
	assume $\divQ(s) = \divQ(s')$, we have (up to a coefficient) $s^{m'} = (s')^m$, which means
	\begin{equation*}
		\lambda / m = \lambda' / m' .
	\end{equation*}
	
	Second, to prove it is a $\mathbb{Q}$-convex linear map, consider two sections
	\begin{equation*}
		s \in H^0(X, m L)^{(B)}_\lambda, \quad s' \in H^0(X, m L)^{(B)}_{\lambda'}
	\end{equation*}
	(it is always possible to find a common $m \in \mathbb{N}(L, \alpha)$) and coefficient $t = {p \over p + q} \in \mathbb{Q} \cap [0, 1]$ (where $p, q \in \mathbb{N}, p + q > 0$). Observe that
	\begin{equation*}
		s^p {s'}^q \in H^0(X, (p + q) m L)^{(B)}_{p \lambda + q \lambda'},
	\end{equation*}
	which results in
	\begin{equation*}
	\begin{aligned}
		p_\alpha(t \cdot \divQ(s) + (1 - t) \cdot \divQ(s')) &= p_\alpha(\divQ(s^p {s'}^q)) \\
		&= {p \lambda + q \lambda' \over (p + q) m} = t \cdot {\lambda \over m} + (1 - t) \cdot {\lambda' \over m} \\
		&= t \cdot p_\alpha(\divQ(s)) + (1 - t) \cdot p_\alpha(\divQ(s')).
	\end{aligned}
	\end{equation*}

	Finally, if $X$ is spherical, the space of $B$-semi-invariant sections $H^0(X, m L)^{(B)}_\lambda$ is $1$-dimensional, so we can define the inverse map 
	\begin{equation*}
		q : P(X, L, \alpha) \to \Qls{L}^B ,
	\end{equation*}
	such that, for every $u = \lambda / m \in {1 \over m} P_m(X, L, \alpha)$, there is a unique (up to scaling) section $s \in H^0(X, m L)^{(B)}_\lambda$. Note that $\divQ(s)$ is independent of the choice of $m$. It is straightforward to check the convex-linearity of the map $q$ by repeating the argument above.
\end{proof}

\begin{example}
	For a flag $G$-variety $X$, and an ample $\mathbb{Q}$-line bundle $L$, the $B$-invariant $\mathbb{Q}$-linear series $\Qls{L}^B$ is a singleton. That means, there is a unique effective $B$-invariant $\mathbb{Q}$-divisor $\mathbb{Q}$-linearly equivalent to $L$.
\end{example}